\subsection{The main results}


The first result concerns the lifting of curves. We~recall that, since $G$ is finite, each continuous $a \colon I \to \si(V)$, where $I\subseteq \R$ is an interval,
has a continuous lift $\ol a \colon I \to V$, by \cite[Theorem 5.1]{LMRac}.



\begin{Theorem}\label{main}
 Let $G$ be a finite group and let $G \acts V$ be a representation of~$G$ on a finite-dimensional
 complex vector space $V$. Let $\si=(\si_1,\dots,\si_n)$ be a (minimal) system of homogeneous
 basic invariants of degrees $d_1,\dots,d_n$ and set $d = \max_i d_i$.
 Let $a \in C^{d-1,1}([\al,\be],\si(V))$ be a~curve
 defined on an open bounded interval
 $(\al,\be)$ with values in $\si(V)$.
 Then each continuous lift $\overline a \colon (\al,\be) \to V$ of $a$ over $\si$ is absolutely continuous and
 belongs to $W^{1,p}((\al,\be),V)$
 with
 \begin{gather} \label{eq:main}
 	 \|\ol a'\|_{L^p((\al,\be))} \le
 C(G \acts V,(\be-\al),p)\, \max_{1 \le j \le n} \|a_j\|^{1/d_j}_{C^{d-1,1}([\al, \be])}		
 	\end{gather}	
 for all $1 \le p < d/(d-1)$,
 where $C$ is a constant which depends only on the representation $G \acts V$, the length of the interval $(\al,\be)$, and $p$.	
\end{Theorem}

{\samepage
The conclusion of the theorem is in general optimal among Sobolev spaces, the differentiability assumption on $a$ is best possible;
see Remark~\ref{optimality}.
Here and below we use the notation
\begin{gather*}
 C^{d-1,1}([\al,\be],\si(V)) := C^{d-1,1}([\al,\be],\C^n) \cap \si(V)^{(\al,\be)},
\end{gather*}
the H\"older class $C^{d-1,1}$ is defined in Section~\ref{functionspaces}.

}
